\documentclass[11pt]{report}
\usepackage{fontspec}
\setmainfont{TeX Gyre Pagella}
\usepackage[margin=1.15in]{geometry}
\usepackage{amsmath,amssymb,amsthm}
\usepackage{enumitem}
\usepackage{hyperref}
\hypersetup{colorlinks=true, linkcolor=black, urlcolor=black, citecolor=black}

\newtheorem{proposition}{Proposition}[chapter]
\newtheorem{theorem}{Theorem}[chapter]
\newtheorem{lemma}{Lemma}[chapter]
\newtheorem{definition}{Definition}[chapter]
\newtheorem{remark}{Remark}[chapter]
\newtheorem{corollary}{Corollary}[chapter]
\newtheorem{example}{Example}[chapter]
\newtheorem{conjecture}{Conjecture}[chapter]

\title{\textbf{Generative Continuation Geometry}\\[0.3em]
\large Working Draft}
\author{Flyxion}
\date{July 2026}

\begin{document}

\maketitle

\tableofcontents

\chapter{Continuation Attractors}

A fern frond is not drawn leaf by leaf. Look closely at one pinna and
it is, near enough, a smaller fern; look closer still at one pinnule
of that pinna and it is, near enough, a smaller fern again. Nobody
specified the shape of the whole plant directly. A small number of
growth rules, applied to their own output over and over, produced
something with more detail than any of the rules themselves contain.
The shape is not stored anywhere. It is generated.

\emph{Continuation Geometry} spent its length measuring futures:
how much room a state has, how far apart two states' futures sit, when
a distinction between futures collapses to nothing. All of that took
the space of futures as given, something to be measured once it
already existed. This book asks a different question. Can a space of
futures itself be grown, the way the fern is grown, from a small
number of simple moves applied to their own output? If continuation
spaces can be generated rather than merely measured, the natural unit
of study stops being the space and becomes the small set of
operations that built it.

This opening chapter answers that question in the one case where the
answer is already, in a precise sense, classical. It is not a new
theorem. It is an old and beautiful one, due to Hutchinson, applied
here to the specific metric geometry the first book already built. The
payoff is that self-similar continuation structures, the kind
suggested by a duck, a crown, and a leaf all arising from the same
underlying rule, turn out to rest on exactly the Hausdorff-distance
machinery already in hand, with nothing new invented and nothing
smuggled in from outside.

\section*{Preliminaries: expansion is not generation}

Two facts from \emph{Continuation Geometry} are worth restating here,
briefly and as what they actually are: lemmas, not theorems, and
useful precisely because they show what generation is \emph{not}.

\begin{lemma}[Reachable expansion]
For $S_t \subseteq X$ and $\mathrm{Adj}(S_t) := \bigcup_{x \in S_t} \Phi(x)$
under a one-step reachability map $\Phi$, the sequence
$S_{t+1} := S_t \cup \mathrm{Adj}(S_t)$ satisfies $S_t \subseteq S_{t+1}$
for every $t$.
\end{lemma}
\begin{proof}
Immediate: $S_t \subseteq S_t \cup \mathrm{Adj}(S_t) = S_{t+1}$ by
definition of union.
\end{proof}

\begin{remark}
This, together with the reachable-set recursion proved by induction
in \emph{Continuation Geometry} (Kauffman's adjacent possible as an
instance of finite-horizon reachability), is the full content of what
that book called generation, and it is worth being honest that it
earns the name only loosely. The sequence $S_0, S_1, S_2, \dots$ grows
monotonically and without bound unless the underlying space is
finite; nothing about it converges, and nothing about it is
self-similar. It is accumulation, not generation in the sense the
fern illustrates. What follows is different in kind: rather than a
sequence of ever-larger unions, it is a fixed operator applied
repeatedly to itself, and the interesting behavior is not growth but
convergence to a shape that reproduces its own generating rule.
\end{remark}

\section*{The space of shapes}

Fix $X = \mathbb{R}^n$ (or any complete metric space) and let
$\mathcal{K}(X)$ denote the set of nonempty compact subsets of $X$,
equipped with the Hausdorff distance $d_H$ of Chapter 6 of
\emph{Continuation Geometry}.

\begin{theorem}[Completeness of the space of shapes]
\label{thm:complete}
If $(X,d)$ is complete, then $(\mathcal{K}(X), d_H)$ is complete.
\end{theorem}
\begin{proof}[Proof sketch]
Given a Cauchy sequence $\{A_n\}$ in $(\mathcal{K}(X), d_H)$, define
$A := \bigcap_{N \geq 1} \overline{\bigcup_{n \geq N} A_n}$. One shows,
using completeness of $X$ and the Cauchy property of $\{A_n\}$ in
$d_H$, that $A$ is nonempty and compact, and that $d_H(A_n, A) \to 0$.
This is a standard but genuinely technical result in its full form,
requiring care with total boundedness and a diagonal argument to
extract convergent selections from each $A_n$; the complete argument
is given in \cite{falconer2003} and \cite{barnsley1988} and is not
reproduced in full here.
\end{proof}

\begin{remark}
This theorem is doing real work that is easy to take for granted.
Without it, nothing built on top of $(\mathcal{K}(X), d_H)$ would be
guaranteed to have a limit at all; a sequence of shapes could get
closer and closer together, in the sense of Hausdorff distance,
without converging to any actual shape. That this cannot happen, that
the space of compact reachable sets is itself complete whenever the
underlying state space is, is the entire reason the construction
below is guaranteed to produce something rather than merely getting
closer to nothing in particular.
\end{remark}

\section*{Contractions and the Hutchinson operator}

\begin{definition}[Contraction]
$T : X \to X$ is a contraction with constant $c \in [0,1)$ if
$d(Tx, Ty) \leq c\, d(x,y)$ for all $x, y \in X$.
\end{definition}

\begin{lemma}[Contractions act contractively on shapes]
\label{lem:setcontract}
If $T$ is a contraction on $X$ with constant $c$, then for
$A, B \in \mathcal{K}(X)$, $d_H(T(A), T(B)) \leq c\, d_H(A,B)$, where
$T(A) := \{Ta : a \in A\}$.
\end{lemma}
\begin{proof}
Fix $a \in A$. Since $B$ is compact, there is $b \in B$ with
$d(a,b) = \operatorname{dist}(a,B) \leq d_H(A,B)$. Then
$d(Ta,Tb) \leq c\, d(a,b) \leq c\, d_H(A,B)$, and since
$Tb \in T(B)$, $\operatorname{dist}(Ta, T(B)) \leq c\, d_H(A,B)$.
Taking the supremum over $a \in A$ gives
$\sup_{a' \in T(A)} \operatorname{dist}(a', T(B)) \leq c\, d_H(A,B)$.
The symmetric argument, exchanging the roles of $A$ and $B$, gives the
same bound the other way, so $d_H(T(A),T(B)) \leq c\, d_H(A,B)$.
\end{proof}

\begin{lemma}[Unions are 1-Lipschitz]
\label{lem:unionlip}
For finite collections $\{A_i\}_{i=1}^n$, $\{B_i\}_{i=1}^n$ of nonempty
compact sets, $d_H\left(\bigcup_i A_i, \bigcup_i B_i\right) \leq
\max_i d_H(A_i, B_i)$.
\end{lemma}
\begin{proof}
Let $m := \max_i d_H(A_i,B_i)$. For $x \in \bigcup_i A_i$, $x \in A_j$
for some $j$, so $\operatorname{dist}(x, B_j) \leq d_H(A_j,B_j) \leq m$,
and since $B_j \subseteq \bigcup_i B_i$,
$\operatorname{dist}(x, \bigcup_i B_i) \leq \operatorname{dist}(x,B_j) \leq m$.
Taking the supremum over $x \in \bigcup_i A_i$ gives
$\sup_{x} \operatorname{dist}(x, \bigcup_i B_i) \leq m$. The symmetric
argument, exchanging the roles of the two collections, gives the same
bound the other way, so
$d_H\left(\bigcup_i A_i, \bigcup_i B_i\right) \leq m$.
\end{proof}

\begin{definition}[Hutchinson operator]
Given contractions $T_1, \dots, T_n$ on $X$ with constants
$c_1, \dots, c_n$, define $\mathcal{H} : \mathcal{K}(X) \to \mathcal{K}(X)$
by $\mathcal{H}(A) := \bigcup_{i=1}^n T_i(A)$.
\end{definition}

\begin{proposition}[$\mathcal{H}$ is a contraction on shapes]
\label{prop:hcontract}
$d_H(\mathcal{H}(A), \mathcal{H}(B)) \leq c\, d_H(A,B)$ for
$c := \max_i c_i < 1$.
\end{proposition}
\begin{proof}
By Lemma~\ref{lem:unionlip},
$d_H(\mathcal{H}(A), \mathcal{H}(B)) = d_H\left(\bigcup_i T_i(A), \bigcup_i T_i(B)\right)
\leq \max_i d_H(T_i(A), T_i(B))$. By Lemma~\ref{lem:setcontract},
$d_H(T_i(A), T_i(B)) \leq c_i\, d_H(A,B)$ for each $i$. Hence
$d_H(\mathcal{H}(A), \mathcal{H}(B)) \leq \max_i c_i\, d_H(A,B) = c\, d_H(A,B)$.
\end{proof}

\section*{The Continuation Attractor Theorem}

\begin{theorem}[Continuation Attractor Theorem, after Hutchinson]
\label{thm:attractor}
Let $T_1, \dots, T_n$ be contractions on a complete metric space $X$
with constants $c_1, \dots, c_n$, and let $c := \max_i c_i < 1$. Then
there exists a unique $A^\ast \in \mathcal{K}(X)$ with
$\mathcal{H}(A^\ast) = A^\ast$, and for any $A_0 \in \mathcal{K}(X)$,
\[
d_H\big(\mathcal{H}^k(A_0), A^\ast\big) \leq \frac{c^k}{1-c}\, d_H(A_0, \mathcal{H}(A_0)) \xrightarrow{k \to \infty} 0.
\]
\end{theorem}
\begin{proof}
By Theorem~\ref{thm:complete}, $(\mathcal{K}(X), d_H)$ is complete. By
Proposition~\ref{prop:hcontract}, $\mathcal{H}$ is a contraction on
this complete metric space. The Banach Fixed Point Theorem then
supplies a unique fixed point $A^\ast$ and the stated geometric
convergence rate, applied directly to $\mathcal{H}$ on
$(\mathcal{K}(X), d_H)$.
\end{proof}

\begin{remark}
No part of this proof required inventing anything: completeness of
$\mathcal{K}(X)$ is Theorem~\ref{thm:complete}, the contraction
property of $\mathcal{H}$ is Proposition~\ref{prop:hcontract} built
from two short lemmas, and the conclusion is the Banach Fixed Point
Theorem, unmodified. This is the sense in which this chapter is
mathematically conservative even though its subject, self-similar
generated shape, sounds exotic: every piece was already sitting in
the metric structure Chapter 6 of \emph{Continuation Geometry}
built for an entirely different purpose.
\end{remark}

\section*{A worked attractor}

\begin{example}[Two refusals]
Let $X = [0,1]$, and let $T_1(x) := x/3$, $T_2(x) := x/3 + 2/3$, both
contractions with constant $c = 1/3$. The attractor $A^\ast$ satisfying
$A^\ast = T_1(A^\ast) \cup T_2(A^\ast)$ is the standard middle-thirds
Cantor set.
\end{example}

\begin{remark}
Read in the vocabulary of continuation, $T_1$ and $T_2$ are the two
survivors of a repeated refusal: at every stage, whatever occupies the
middle third of the current set of admissible outcomes is discarded,
and the remaining two-thirds are each rescaled to stand in for the
whole. What survives in the limit is neither of the two pieces nor
their simple union at any finite stage, but the unique shape that
reproduces itself exactly under the same rule that built it,
uncountable, nowhere dense, and entirely determined by two contraction
maps. Whether operators named Collapse or Refuse elsewhere in this
corpus are literally instances of $T_i$ in this technical sense, maps
satisfying the contraction inequality on some genuine underlying
metric, is a claim about those specific constructions and is not
asserted here; what is established is only that \emph{if} a family of
continuation-narrowing operators can be shown to be contractions on
$(\mathcal{K}(X), d_H)$, this theorem applies to them without
modification.
\end{remark}

\section*{What this chapter does not yet claim}

This chapter establishes that self-similar continuation structures are
mathematically available, not that any particular structure this
corpus has described informally is an instance of one. Two gaps in
particular are worth naming rather than leaving implicit. First,
nothing here says which operators arising from an actual continuation
problem, as opposed to a hand-picked pair of affine maps on $[0,1]$,
are genuinely contractions on a genuine metric; that has to be checked
case by case, the way it would for any application of a fixed-point
theorem. Second, and more speculatively, the question of whether a
small, fixed family of primitive operators could generate \emph{every}
admissible continuation geometry to arbitrary Hausdorff-distance
accuracy, a continuation-geometric analogue of universal approximation,
is not addressed by this theorem and is not claimed here. It is a
genuine open question, and answering it is a natural place for this
book to go next, once the machinery in the chapters between here and
there, quotient maps and operator composition among them, has been
built on equally solid ground.

\chapter{Continuation Quotients}

Two rivers can rise on opposite sides of a mountain range, cross
entirely different terrain, pass through valleys that share nothing
in common, and still arrive at the same confluence and become, from
that point on, one river. Everything about their separate histories,
which rocks each one wore smooth, which tributaries fed it, how long
each took to descend, becomes irrelevant downstream of the merge. A
boat launched below the confluence cannot tell which branch its water
came from. Not because the information was hidden. Because it no
longer makes a difference to anything the boat will experience from
here on.

The previous chapter showed that a continuation space can be grown
from a small number of generating operators. The natural next
question is what happens when two different generating histories,
different sequences of operators, different starting points, produce
the same continuation space at the end. When does that happen, and
what, if anything, does it mean? The temptation is to reach
immediately for language like information loss or memory collapse,
language already used loosely elsewhere in this corpus. This chapter
resists that temptation on purpose. It defines the relevant object
first, proves what can be proved about it without reference to any
interpretation, and only afterward asks what, if anything, has been
lost.

\section*{The continuation quotient}

\begin{definition}[Continuation quotient]
Let $H$ be a set of histories and let $\Gamma : H \to \mathcal{K}(X)$
assign to each history its continuation space (or, following Chapter 6
of \emph{Continuation Geometry}, its finite-horizon reachable set
$R_T$ for a fixed horizon $T$). Define $h_1 \sim_\Gamma h_2$ iff
$\Gamma(h_1) = \Gamma(h_2)$. This is an equivalence relation
(reflexivity, symmetry, and transitivity all following immediately
from equality of sets), and the continuation quotient is
$H/{\sim_\Gamma}$, with quotient map $q : H \to H/{\sim_\Gamma}$
sending $h$ to its equivalence class $[h]$.
\end{definition}

Nothing has been said yet about distinctions, memory, or loss. This is
deliberate: the object exists, and has structure, independently of any
of those interpretations.

\begin{proposition}[Well-defined continuation projection]
\label{prop:universal}
There exists a unique map $\bar\Gamma : H/{\sim_\Gamma} \to \mathcal{K}(X)$
such that $\Gamma = \bar\Gamma \circ q$.
\end{proposition}
\begin{proof}
Define $\bar\Gamma([h]) := \Gamma(h)$. This is well-defined: if
$[h_1] = [h_2]$, then $h_1 \sim_\Gamma h_2$, so
$\Gamma(h_1) = \Gamma(h_2)$ by definition of $\sim_\Gamma$, so the
value assigned to $[h_1] = [h_2]$ does not depend on which
representative was chosen. By construction,
$\bar\Gamma(q(h)) = \bar\Gamma([h]) = \Gamma(h)$ for every $h \in H$,
so $\Gamma = \bar\Gamma \circ q$. For uniqueness, suppose
$\bar\Gamma'$ also satisfies $\Gamma = \bar\Gamma' \circ q$. Since $q$
is surjective, every element of $H/{\sim_\Gamma}$ is $q(h)$ for some
$h$, and $\bar\Gamma'(q(h)) = \Gamma(h) = \bar\Gamma(q(h))$ for every
such $h$, so $\bar\Gamma' = \bar\Gamma$.
\end{proof}

\begin{remark}
This is routine, and it is worth saying plainly that it is routine:
it is the universal property of a quotient set, applied to nothing
more exotic than the map $\Gamma$ itself. What makes it worth stating
is not its difficulty but what it licenses: from here on,
$\bar\Gamma$ is a genuine function on the quotient, not merely a
suggestive notation, and anything proved about $\bar\Gamma$ is a fact
about the quotient set itself, not about any particular history used
to reach a point in it.
\end{remark}

\section*{A structural lemma}

The worked example below relies on the following fact about
time-homogeneous, deterministic-image dynamics, worth isolating here.

\begin{lemma}[One-step agreement propagates]
\label{lem:propagate}
Suppose $R_1(x) = R_1(y)$ for states $x, y$ under a time-homogeneous
reachability map. Then $R_T(x) = R_T(y)$ for every $T \geq 1$.
\end{lemma}
\begin{proof}
By induction on $T$. The case $T=1$ is the hypothesis. Suppose
$R_T(x) = R_T(y)$. Since the dynamics are time-homogeneous,
$R_{T+1}(x) = \bigcup_{z \in R_T(x)} R_1(z)$ and likewise for $y$; since
$R_T(x) = R_T(y)$, these unions range over the same set of states $z$,
so $R_{T+1}(x) = R_{T+1}(y)$.
\end{proof}

\begin{remark}
This is the precise sense in which a collision, once it happens, is
permanent under time-homogeneous dynamics: two states whose immediate
futures coincide exactly will have every subsequent future coincide
as well, since there is nothing left in either state's future to
distinguish it from the other's.
\end{remark}

\section*{Quotient invariants}

\begin{proposition}[Quotient invariance]
\label{prop:invariance}
Let $F : \mathcal{K}(X) \to Y$ be any function. Then
$F \circ \Gamma : H \to Y$ factors through $q$: there is a unique
$\bar{F} : H/{\sim_\Gamma} \to Y$ with $F \circ \Gamma = \bar{F} \circ q$.
\end{proposition}
\begin{proof}
$F \circ \Gamma : H \to Y$ satisfies $(F \circ \Gamma)(h_1) = (F \circ \Gamma)(h_2)$
whenever $h_1 \sim_\Gamma h_2$, since $\Gamma(h_1) = \Gamma(h_2)$ by
definition of $\sim_\Gamma$ and $F$ is a function, hence single-valued.
Apply Proposition~\ref{prop:universal} with $F \circ \Gamma$ in place
of $\Gamma$ to obtain the factorization.
\end{proof}

\begin{corollary}
Continuation volume $\operatorname{Vol}_T$, the distinction margin
$A_{\mathcal D}^T$ relative to any fixed reference distinction, and
boundary proximity $\delta$, all defined in \emph{Continuation
Geometry} as functions of a reachable set or of a state's position
relative to one, descend to well-defined functions on
$H/{\sim_\Gamma}$.
\end{corollary}
\begin{proof}
Each is of the form $F(\Gamma(h))$ or $F(R_T(x(h)))$ for a function
$F$ depending only on the reachable set itself (Lebesgue measure of
the set, Hausdorff distance to another fixed set, ambient distance
from a fixed reference set), so Proposition~\ref{prop:invariance}
applies directly to each.
\end{proof}

\begin{remark}
This is a broader statement than it might first appear. It says that
any quantity this book or its predecessor has defined, or will define,
purely in terms of a reachable set, without reference to which
history produced that set, is automatically constant on
$\sim_\Gamma$-equivalence classes, with no further argument required
each time a new such quantity is introduced. The work of checking
invariance, done once here, does not need to be redone for every
future functional built the same way.
\end{remark}

\section*{Finite checkability}

Lemma~\ref{lem:propagate} shows that equal one-step images propagate
forward forever. It says nothing about the converse: two states whose
one-step images differ can still coincide at some later horizon.
Nothing in this chapter has claimed otherwise, and the general
question, whether continuation equivalence at every horizon can still
be checked in finitely many steps despite this, deserves a correct
answer rather than an implied one.

\begin{proposition}[Eventual periodicity bounds the check]
\label{prop:periodicity}
If $X$ is finite, then for any $x \in X$ the sequence
$R_1(x), R_2(x), R_3(x), \dots$ is eventually periodic, with both the
pre-period and the period bounded by $2^{|X|}$. Consequently, whether
$R_T(x) = R_T(y)$ holds for all sufficiently large $T$ can be
determined by comparing $R_T(x)$ and $R_T(y)$ for $T$ ranging only up
to $2 \cdot 2^{|X|}$.
\end{proposition}
\begin{proof}
Each $R_T(x)$ is an element of $2^X$, which has $2^{|X|}$ elements, so
the sequence $R_1(x), R_2(x), \dots$ takes at most $2^{|X|}$ distinct
values; among its first $2^{|X|}+1$ terms, two must coincide by
pigeonhole, $R_p(x) = R_q(x)$ for some $1 \leq p < q \leq 2^{|X|}+1$.
Time-homogeneity means $R_{T+1} = \Psi(R_T)$ for the fixed map
$\Psi(S) := \bigcup_{z \in S} \Phi(z)$, so $R_p(x) = R_q(x)$ gives
$R_{p+j}(x) = \Psi^j(R_p(x)) = \Psi^j(R_q(x)) = R_{q+j}(x)$ for every
$j \geq 0$: the sequence is periodic with period dividing $q - p \leq
2^{|X|}$ from index $p \leq 2^{|X|}$ onward. The same bound applies to
the sequence generated by $y$. Once both sequences have entered their
eventual cycles, comparing them over one further full period, at most
$2^{|X|}$ more steps, settles whether they agree from that point on.
\end{proof}

\begin{remark}
This is the correct version of the claim that continuation
equivalence is finitely checkable. It is a coarser bound than a sharp
analysis of any particular system would give, since $2^{|X|}$ is a
worst case over all possible dynamics on a state space of that size,
but it is unconditional: no structural assumption on $\Phi$ beyond
finiteness of $X$ and time-homogeneity is needed. The traffic light
below reaches its own eventual behavior in four steps, far inside this
bound, and it is worth seeing exactly how.
\end{remark}

\section*{The traffic light, revisited}

Recall the traffic light of \emph{Continuation Geometry}, Chapter 10:
states $(k, r) \in \{0,1,2,3\} \times \{0,1\}$, phase $k$ and pending
request $r$. To compute $\Gamma$ properly requires treating the
system's branching explicitly rather than leaving it implicit. Let
$e \in \{0,1\}$ record whether a new pedestrian request arrives during
a given step (exogenous, not controlled by the phase). The transition
is
\[
\Phi((k,r), e) = \left( (k+1) \bmod 4,\ r' \right), \qquad
r' = \begin{cases}
0 & k+1 \equiv 0 \pmod 4 \text{ (clear on arrival at red)} \\
1 & k+1 \not\equiv 0 \pmod 4 \text{ and } (r = 1 \text{ or } e = 1) \\
0 & k+1 \not\equiv 0 \pmod 4 \text{ and } r = 0 \text{ and } e = 0.
\end{cases}
\]
Let $H$ be the set of finite paths of the system, and
$\Gamma(h) := R_1(x(h))$, the one-step reachable set from the
endpoint of $h$, ranging over $e \in \{0,1\}$.

\begin{proposition}[Exact computation of the traffic-light quotient at $T=1$]
\label{prop:tlquotient}
Under $\sim_\Gamma$, the eight states fall into exactly seven classes:
the single nontrivial class $\{(3,0),(3,1)\}$, and six singleton
classes, one for each remaining state.
\end{proposition}
\begin{proof}
Compute $R_1$ for all eight states directly from $\Phi$. For
$k \neq 3$, arrival is at a nonzero phase and the ordinary branching
rule applies, $r' = 1$ iff $r=1$ or $e=1$:
\begin{align*}
R_1(0,0) &= \{(1,0),(1,1)\}, & R_1(0,1) &= \{(1,1)\}, \\
R_1(1,0) &= \{(2,0),(2,1)\}, & R_1(1,1) &= \{(2,1)\}, \\
R_1(2,0) &= \{(3,0),(3,1)\}, & R_1(2,1) &= \{(3,1)\}.
\end{align*}
For $k=3$, arrival is at phase $0$, so clearing fires unconditionally
regardless of $r$ or $e$:
\[
R_1(3,0) = \{(0,0)\}, \qquad R_1(3,1) = \{(0,0)\}.
\]
Comparing all eight sets directly: the three size-two sets,
$\{(1,0),(1,1)\}$, $\{(2,0),(2,1)\}$, $\{(3,0),(3,1)\}$, are pairwise
distinct; the four remaining size-one sets, $\{(1,1)\}$, $\{(2,1)\}$,
$\{(0,0)\}$ (appearing twice, from $(3,0)$ and $(3,1)$), and
$\{(0,0)\}$ again, give exactly one coincidence, $R_1(3,0) = R_1(3,1)
= \{(0,0)\}$, and no other.
\end{proof}

Now define the mechanical quotient $\sim_M$ independently of
$\Gamma$, as the identification the design actually performs: the
clearing rule fires on the transition \emph{into} phase $0$, so it is
the two states about to make that transition, $(3,0)$ and $(3,1)$,
that the design treats as interchangeable, both being followed
unconditionally by $(0,0)$.

\begin{proposition}[The design is exact at $T=1$]
$\sim_M\ =\ \sim_\Gamma$.
\end{proposition}
\begin{proof}
Both relations have exactly one nontrivial class, $\{(3,0),(3,1)\}$,
and agree as singletons on every other state: $\sim_M$ by definition
of the clearing rule, $\sim_\Gamma$ by Proposition~\ref{prop:tlquotient}.
Two equivalence relations on the same finite set with identical
partitions are equal.
\end{proof}

\begin{remark}
It is worth being precise about which of two similar-sounding designs
this result actually vindicates. Clearing on arrival at phase $0$
(the rule above) is exact. A superficially similar alternative,
clearing whenever the system is currently \emph{at} phase $0$ rather
than arriving there, would instead identify $(0,0)$ and $(0,1)$. But
$R_1(0,0) = \{(1,0),(1,1)\} \neq \{(1,1)\} = R_1(0,1)$: a state at the
red light with no pending request can still branch toward either
outcome next, while one with a pending request cannot, precisely
because that request is about to be served. The alternative design
would merge two states with genuinely different one-step futures,
exactly the coarser, information-destroying case this chapter opened
by distinguishing from the exact one. The two designs differ only in
whether clearing is keyed to the transition or to the state, and only
one of them is correct.
\end{remark}

\section*{What one horizon hides}

Restricting Proposition~\ref{prop:tlquotient} to $T=1$ was not a
simplification; it is what $\Gamma$ was defined to measure. But
Proposition~\ref{prop:periodicity} guarantees the longer-horizon
question has a definite, finitely checkable answer, and for this
system that answer is worth computing directly rather than left as an
abstract guarantee.

\begin{proposition}[The request bit is eventually redundant everywhere]
\label{prop:eventual}
For each phase $k$, the pair $(k,0)$ and $(k,1)$ satisfies
$R_T(k,0) = R_T(k,1)$ for all $T \geq 4 - k$, and not before.
\end{proposition}
\begin{proof}[Proof by direct computation]
By Proposition~\ref{prop:tlquotient}, $(3,0) \sim_\Gamma (3,1)$ at
$T=1 = 4-3$. Applying $\Psi$ once more,
$R_2(2,0) = \Psi(R_1(2,0)) = \Psi(\{(3,0),(3,1)\})$ and
$R_2(2,1) = \Psi(R_1(2,1)) = \Psi(\{(3,1)\})$; since
$\Phi(3,0,e) = \Phi(3,1,e) = (0,0)$ for both $e$, both unions reduce
to $\{(0,0)\}$, so $R_2(2,0) = R_2(2,1) = \{(0,0)\}$, matching
$T = 2 = 4-2$. The same one-step displacement argument gives
$R_3(1,0) = R_3(1,1)$ at $T=3=4-1$ and $R_4(0,0) = R_4(0,1)$ at
$T=4=4-0$: each pair's images are, one step earlier, exactly the
pair merged at the previous phase, so equality propagates forward by
exactly one horizon per phase step. That equality does not hold
earlier follows from the cardinality mismatch already established at
$T=1$ in Proposition~\ref{prop:tlquotient}, propagated the same one
step at a time: at $T = 4-k-1$, the images of $(k,0)$ and $(k,1)$ are
the not-yet-merged pair one phase ahead.
\end{proof}

\begin{remark}
This is the sharper and more interesting fact the single-horizon
result was hiding. The request bit is not redundant in general; it is
redundant \emph{by a deadline}, and that deadline is exactly the
number of steps remaining until the light is guaranteed to cycle
through its clearing transition, four steps for a state that has just
left it, one step for a state about to enter it. Every state's
request bit stops mattering to its continuation space at precisely
the moment continuation stops being able to tell whether that request
will still be pending when clearing finally arrives, and not one step
before. The design clears the bit at the earliest possible moment,
$T=1$ after arrival, which is also the latest moment at which failing
to clear it would still have been observably wrong; every other
phase's bit becomes redundant on its own, without any design
intervention, simply because the light's own periodicity eventually
answers the question the bit was tracking.
\end{remark}

\section*{What quotients do not yet explain}

This chapter has shown when two histories are interchangeable from
the standpoint of every quantity built purely from continuation
spaces, and it has shown, in one worked case, that a designed
identification can be checked against that standard exactly rather
than merely defended by intuition. It has not yet said anything about
what is lost when a quotient is coarser than $\sim_\Gamma$, only that
the traffic light's particular quotient happens not to be. The next
chapter takes up the coarser case directly: not whether information is
destroyed, which this chapter's machinery can already detect, but how
much, and by what measure.

\chapter{Distinction Loss Under Quotients}

A library catalog entry is a quotient of a book: title, author, and a
handful of subject headings standing in for everything between the
covers. A good catalog loses nothing a patron actually needs in order
to decide whether to pull the book from the shelf. A bad one can lose
quite a lot; a catalog that records only ``fiction'' or
``nonfiction'' will happily file a memoir next to a novel and a field
guide next to a textbook, and a patron relying on that heading alone
has no way to tell them apart until the book is already in hand. The
catalog entry is not wrong, exactly; it simply was not built to
preserve the distinctions this particular patron cares about.

The previous chapter showed how to check whether a given
identification of histories is exact: whether it merges only
histories whose continuation spaces already coincide, or whether it
merges some that do not. That check was binary, exact or not. This
chapter asks the sharper question a librarian actually has to answer.
When an identification is not exact, how much is lost, and where
exactly does the loss occur? The traffic light's badly designed
alternative from the previous chapter turns out to have a precise
answer, not just a qualitative one.

\section*{Fibers and their content}

\begin{definition}[Fiber, $\Gamma$-content, loss]
For a quotient $q : H \to O$ and $o \in O$, the fiber over $o$ is
$q^{-1}(o) \subseteq H$. Its $\Gamma$-content is
$C_\Gamma(o) := \{\Gamma(h) : h \in q^{-1}(o)\}$, the set of distinct
continuation spaces appearing in the fiber. The loss at $o$ is
\[
\mathcal{L}(o) := \log |C_\Gamma(o)|,
\]
so $\mathcal{L}(o) = 0$ exactly when every history mapping to $o$
under $q$ shares a single continuation space.
\end{definition}

\begin{proposition}[Zero loss characterizes $\Gamma$-respecting quotients]
\label{prop:zeroloss}
$\mathcal{L}(o) = 0$ for every $o \in O$ if and only if $q$ refines
$\sim_\Gamma$: $h_1 \sim_q h_2 \implies h_1 \sim_\Gamma h_2$.
\end{proposition}
\begin{proof}
If $q$ refines $\sim_\Gamma$, then $h_1 \sim_q h_2$ implies
$\Gamma(h_1) = \Gamma(h_2)$, so every fiber of $q$ lies inside a
single $\sim_\Gamma$-class and has $|C_\Gamma(o)| = 1$, giving
$\mathcal{L}(o) = 0$. Conversely, if $\mathcal{L}(o) = 0$ for every
$o$, then every fiber has $|C_\Gamma(o)| = 1$, so any two histories in
the same fiber share the same $\Gamma$-value: $h_1 \sim_q h_2 \implies
\Gamma(h_1) = \Gamma(h_2) \implies h_1 \sim_\Gamma h_2$, so $q$
refines $\sim_\Gamma$.
\end{proof}

\section*{The coarsest lossless quotient}

\begin{theorem}[$\sim_\Gamma$ is the coarsest zero-loss quotient]
\label{thm:coarsest}
Every zero-loss quotient $q : H \to O$ refines $\sim_\Gamma$, and
$\sim_\Gamma$ itself has zero loss everywhere. Consequently
$\sim_\Gamma$ is, up to relabeling of $O$, the unique coarsest
quotient of $H$ with zero loss.
\end{theorem}
\begin{proof}
The first claim is Proposition~\ref{prop:zeroloss}. That
$\sim_\Gamma$ itself has zero loss is immediate: by definition,
$h_1 \sim_\Gamma h_2$ already means $\Gamma(h_1) = \Gamma(h_2)$, so
every $\sim_\Gamma$-class has exactly one $\Gamma$-value. Since every
zero-loss $q$ refines $\sim_\Gamma$, no zero-loss quotient can be
strictly coarser than $\sim_\Gamma$, and $\sim_\Gamma$ achieves zero
loss, so it is the coarsest one.
\end{proof}

\begin{remark}
This is the quotient-theoretic form of a fact Chapter 5 of
\emph{Continuation Geometry} established for sufficient statistics
without this vocabulary. A sufficient statistic $\varphi$ there was
required to satisfy $\varphi(H_t) = \varphi(H_t') \implies
\Gamma(x_t \mid H_t) = \Gamma(x_t \mid H_t')$, which is exactly the
condition that the quotient induced by $\varphi$ refines
$\sim_\Gamma$. Theorem~\ref{thm:coarsest} says $\sim_\Gamma$ itself is
the coarsest possible sufficient statistic: the most compressed
summary of a history that still loses nothing continuation-relevant.
Any coarser summary is not merely less detailed; by
Proposition~\ref{prop:zeroloss}, it has strictly positive loss
somewhere.
\end{remark}

\begin{proposition}[Coarsening never decreases loss]
\label{prop:monotone}
If $q'$ refines $q$ (every $q'$-fiber is contained in some $q$-fiber),
then for every $q$-fiber $q^{-1}(o)$, expressed as a union of
$q'$-fibers $q'^{-1}(o'_1), \dots, q'^{-1}(o'_k)$,
\[
\mathcal{L}_q(o) \geq \max_i \mathcal{L}_{q'}(o'_i).
\]
\end{proposition}
\begin{proof}
$C_\Gamma(o) = \bigcup_i C_{\Gamma}(o'_i)$, since the $\Gamma$-values
appearing in $q^{-1}(o)$ are exactly those appearing in its
constituent $q'$-fibers. A union of sets is at least as large as any
one of them, so $|C_\Gamma(o)| \geq |C_\Gamma(o'_i)|$ for each $i$,
and $\log$ is monotone increasing.
\end{proof}

\begin{remark}
Merging two fibers can only add to the set of continuation spaces an
observer sees when looking at the merged fiber; it can never remove
one. This is the precise sense in which coarsening a quotient is a
one-way trip toward loss: refining a lossy quotient can repair the
loss, by splitting a mixed fiber back into pieces that each see only
one $\Gamma$-value, but no further coarsening of an already-lossy
quotient can undo the damage, only add to it or leave it unchanged.
\end{remark}

\section*{Distinction margin as a local loss detector}

\begin{proposition}[Vanishing margin is fiber collision]
For a distinction $\mathcal{D}$ and states $x, y$ in different cells
of $\mathcal{D}$, $A_{\mathcal D}^T(x) = 0$ (Chapter 7 of
\emph{Continuation Geometry}) if and only if $x$ and $y$ lie in the
same fiber of the quotient induced by $\mathcal D$ but in different
$\sim_\Gamma$-classes at horizon $T$: that is, $\mathcal D$ merges $x$
and $y$ while $R_T(x) \neq R_T(y)$ for some witnessing $y$, or
$R_T(x) = R_T(y)$ for the specific $y$ achieving the infimum.
\end{proposition}
\begin{proof}
Recall $A_{\mathcal D}^T(x) := \inf_{y \notin D(x)} d_H(R_T(x), R_T(y))$.
This infimum is $0$ exactly when some $y$ outside $x$'s
$\mathcal{D}$-cell has $R_T(x) = R_T(y)$ (attained, under the
compactness hypothesis already used in Chapter 7), which is exactly
the statement that $x$ and that $y$ share a $\Gamma$-value at horizon
$T$ despite $\mathcal D$ having placed them in different cells: they
are $\sim_\Gamma$-equivalent (at horizon $T$) even though $\mathcal D$
distinguishes them.
\end{proof}

\begin{remark}
This gives the distinction margin of the first book a second reading.
It was introduced there as a measure of how close a state sits to
losing a distinction. Read through this chapter's vocabulary, it is
something more specific: a pointwise diagnostic for whether the
observer's chosen distinction $\mathcal{D}$, treated as its own
quotient of the state space, disagrees with $\sim_\Gamma$ at that
particular point. $A_{\mathcal D}^T(x) = 0$ does not merely mean $x$
is close to indistinguishable from something outside its cell; by
this chapter's Proposition~\ref{prop:zeroloss} applied locally, it
means $\mathcal D$'s quotient has positive loss exactly there,
whether or not it has loss anywhere else.
\end{remark}

\section*{The traffic light, quantified}

Return to the two designs compared in the previous chapter: the
correct one, clearing on arrival at phase $0$, and the alternative,
clearing whenever currently at phase $0$.

\begin{example}[Zero loss, and one bit]
The correct design's quotient is $\sim_\Gamma$ itself (Proposition~4
of the previous chapter), so by Theorem~\ref{thm:coarsest} its loss is
$\mathcal{L}(o) = 0$ at every fiber $o$. The alternative design merges
$(0,0)$ and $(0,1)$ into a single fiber $o_{\mathrm{bad}}$. Their
$\Gamma$-values at $T=1$ are $R_1(0,0) = \{(1,0),(1,1)\}$ and
$R_1(0,1) = \{(1,1)\}$, which are distinct, so
$C_\Gamma(o_{\mathrm{bad}}) = \{\{(1,0),(1,1)\},\ \{(1,1)\}\}$, a set
of size $2$, giving
\[
\mathcal{L}(o_{\mathrm{bad}}) = \log 2,
\]
exactly one bit, and $\mathcal{L}(o) = 0$ at every other fiber, since
the alternative design agrees with the correct one everywhere except
at this single merge.
\end{example}

\begin{remark}
This is what the previous chapter's qualitative verdict, that the
alternative design was coarser and information-destroying, actually
amounts to once measured: not a diffuse degradation spread across the
whole system, but a single, exactly located bit of loss, sitting at
precisely the fiber where the alternative design merged two states
that a patron of this particular library, a controller trying to
plan four steps ahead, would still have wanted to tell apart. Every
other fiber of the alternative design is, by coincidence or by the
limited scope of its one design error, exactly as good as the correct
one.
\end{remark}

\section*{What loss does not yet explain}

This chapter can say, precisely, how much is lost by a given
quotient and exactly where. It has nothing to say about why a
particular coarser quotient might be chosen anyway, in full knowledge
of the loss it incurs. A designer might prefer a one-bit-lossy
traffic light if the saved bit of memory matters more than the
distinction it would have preserved; a librarian might prefer the
cheap two-category catalog if most patrons never need the finer one.
Measuring loss and deciding whether a loss is worth paying for are
different questions, and this chapter has only ever answered the
first of them.

\chapter{Continuation Operators}

The first three chapters of this book have all, in their own way, been
about nouns: a shape a system of contractions converges to, a
partition of histories that share a future, a number measuring what a
partition throws away. None of them named the verb. A contraction
does something to a shape. A design decision, clearing a bit on
arrival at a red light rather than while sitting at one, does
something to a history. This chapter is about the doing: what counts
as a legitimate move in this framework, and what survives when two
different moves are composed into one.

The two kinds of move already used in this book look, on the surface,
like opposites. Chapter 1's contractions built detail up, each
application adding a smaller, self-similar copy of the whole. Chapters
2 and 3's quotients tore detail down, collapsing histories that had
become, for every purpose this book has formalized, the same. It
would be a missed opportunity to leave them looking like two
unrelated tools. They are both instances of one idea, and this
chapter's job is to say what that idea is.

\section*{Operators and their composition}

\begin{definition}[Continuation operator]
A continuation operator on $X$ is a Lipschitz map
$T : \mathcal{K}(X) \to \mathcal{K}(X)$: there is a constant
$L \geq 0$ with $d_H(T(A), T(B)) \leq L\, d_H(A,B)$ for all
$A, B \in \mathcal{K}(X)$. Write $L(T)$ for the smallest such constant.
\end{definition}

\begin{proposition}[Closure under composition]
\label{prop:closure}
If $T_1, T_2$ are continuation operators, so is $T_2 \circ T_1$, with
$L(T_2 \circ T_1) \leq L(T_2)\, L(T_1)$.
\end{proposition}
\begin{proof}
$d_H(T_2(T_1(A)), T_2(T_1(B))) \leq L(T_2)\, d_H(T_1(A), T_1(B)) \leq
L(T_2)\, L(T_1)\, d_H(A,B)$.
\end{proof}

\begin{remark}
This proof is one line because it has to be; the definition was
chosen precisely so that composition would cost nothing further to
verify. What makes the definition worth having is not this
proposition but what it is a common home for. The contractions of
Chapter 1 are continuation operators with $L(T) < 1$. The map
$\Psi(A) := \bigcup_{z \in A} \Phi(z)$ used throughout Chapters 2 and
3 to advance a reachable set by one step is a continuation operator
too, generally with $L(\Psi) \geq 1$: it need not contract, but it is
still Lipschitz whenever the underlying one-step map $\Phi$ is well
behaved. Both belong to the same category. What distinguishes a
generative operator from a reductive one is not the definition above;
it is what happens under repeated composition, contraction toward a
single attractor in one case, and, as the next section makes precise,
collapse of distinctions in the other.
\end{remark}

\section*{Operators on histories, and what they respect}

The operators above act on shapes directly. A different and equally
natural kind of operator acts on histories: appending one more real
transition, editing a codebase, applying a repair. The interesting
question for an operator of this kind is not its Lipschitz constant
but whether it interacts predictably with $\sim_\Gamma$.

\begin{definition}[$\Gamma$-compatible operator]
$T : H \to H$ is $\Gamma$-compatible if there is a map
$\hat{T} : \mathcal{K}(X) \to \mathcal{K}(X)$ with
$\Gamma(T(h)) = \hat{T}(\Gamma(h))$ for every $h \in H$: the
continuation space of the transformed history depends on the original
history only through its continuation space, not otherwise.
\end{definition}

\begin{proposition}[$\Gamma$-compatible operators descend to the quotient]
\label{prop:descend}
If $T$ is $\Gamma$-compatible, then $h_1 \sim_\Gamma h_2 \implies
T(h_1) \sim_\Gamma T(h_2)$: $T$ induces a well-defined map
$\bar{T} : H/{\sim_\Gamma} \to H/{\sim_\Gamma}$.
\end{proposition}
\begin{proof}
This is Proposition~4 of the previous chapter (quotient invariance),
applied to $F := \hat{T} \circ \Gamma$ in place of $\Gamma$ itself:
$h_1 \sim_\Gamma h_2$ gives $\Gamma(h_1) = \Gamma(h_2)$, hence
$\Gamma(T(h_1)) = \hat{T}(\Gamma(h_1)) = \hat{T}(\Gamma(h_2)) =
\Gamma(T(h_2))$, so $T(h_1) \sim_\Gamma T(h_2)$.
\end{proof}

\begin{remark}
This is the sense in which $\Gamma$-compatibility is the right
condition to ask for: it says an operator cannot secretly depend on
information that continuation equivalence has already declared
irrelevant. An operator failing this condition would be able to treat
two continuation-equivalent histories differently for reasons that
have nothing to do with what either history can still become, which
is exactly the kind of dependence Chapter 5 of \emph{Continuation
Geometry} worried about under the name of history-dependence in the
first place.
\end{remark}

\begin{example}[The traffic light's step operators]
For $e \in \{0,1\}$, let $T_e : H \to H$ append the transition
$\Phi(\cdot, e)$ to a history. Whether $T_e$ is $\Gamma$-compatible
(at horizon $T=1$) needs checking only on the one nontrivial
$\sim_\Gamma$-class found in Chapter 2, $\{(3,0),(3,1)\}$, since every
other class is a singleton and compatibility there is automatic.
Direct computation gives $T_0(3,0) = T_0(3,1) = (0,0)$ and
$T_1(3,0) = T_1(3,1) = (0,0)$: both step operators send the two
equivalent histories to the identical resulting state, not merely to
states with equal $R_1$, so compatibility holds here in the strongest
possible form.
\end{example}

\begin{remark}
This is a narrow check, confined to the single class where there was
anything to check, and it should not be read as a general fact about
$T_0$ or $T_1$ on this system, only as a confirmation that neither
step operator introduces a new dependence on which of $(3,0)$ or
$(3,1)$ a history happened to pass through, at exactly the point
where such a dependence would have mattered.
\end{remark}

\section*{The family $\Gamma_T$ and the operator $\Psi$}

A different relationship, easy to conflate with $\Gamma$-compatibility
but distinct from it, has been used without a name since Chapter 2. It
concerns not an operator on histories but the relationship between the
finite-horizon continuation spaces at consecutive horizons for one and
the same history.

\begin{remark}[$\Psi$ relates horizons, not histories]
By construction, $R_{T+1}(x) = \bigcup_{z \in R_T(x)} \Phi(z) =
\Psi(R_T(x))$, i.e.\ $\Gamma_{T+1} = \Psi \circ \Gamma_T$ as an
identity between functions $H \to \mathcal{K}(X)$, where
$\Gamma_T(h) := R_T(x(h))$. This is not an instance of
Proposition~\ref{prop:descend}; no operator on $H$ is involved, only
the fixed continuation operator $\Psi$ acting on the shape side, once
for each increment of the horizon. It is this identity, not any
property of an operator on histories, that did the real work in
Lemma~2 of the previous chapter (one-step agreement propagates) and in
the eventual-periodicity argument of Chapter 3: both are, in
retrospect, statements about the orbit of a single reachable set
under repeated application of the one continuation operator $\Psi$,
dressed in the language of horizons rather than the language of
operators used here.
\end{remark}

\section*{Two families, one category}

\begin{corollary}
Both the contractions of Chapter 1 and the operator $\Psi$ of Chapters
2 and 3 are continuation operators in the sense of this chapter's
first definition; they differ only in whether their Lipschitz
constant is below $1$.
\end{corollary}

\begin{remark}
This is worth stating plainly because it says the split this book
opened with, generative operators building attractors versus
reductive operators collapsing distinctions, is not a split in what
kind of mathematical object is involved. It is a split in one number.
A continuation operator with $L(T) < 1$ pulls every pair of shapes
closer together under repeated application, which is exactly
Chapter 1's mechanism for generating a unique self-similar attractor.
A continuation operator with $L(T) \geq 1$ need not pull anything
together at all, and $\Psi$ in particular does something else
entirely on a finite state space: not converge to a single shape, but
cycle, as the traffic light's own eventual periodicity already showed.
Generation and collapse are not two theories. They are two regions of
the same parameter.
\end{remark}

\section*{What operators do not yet explain}

This chapter has said what a legitimate move is and shown that both
of this book's central constructions are instances of it, but it has
not asked which finite collections of operators are rich enough to
produce every continuation geometry worth having, only that individual
operators compose predictably once they exist. That question, closer
to the universal-generation intuition this project began with, needs
the vocabulary built here as a prerequisite, and is not attempted in
this chapter.

\chapter{Toward Universal Generation}

A handful of letters generates every word in a dictionary that has
not been written yet. A handful of chess rules generates every game
that will ever be played under them, including games no one has
conceived of. A small, fixed vocabulary reaching an unbounded space
of outcomes is one of the oldest tricks available to any generative
system, and it is worth asking, now that this book has a settled
vocabulary for what a continuation operator is, whether continuation
geometries have the same property. Can a small, fixed family of
operators, chosen once and never enlarged, get arbitrarily close to
any admissible continuation geometry at all?

This chapter does not answer that question. It draws its boundary
precisely enough to say exactly what would need to be true for the
answer to be yes, proves the one direction that is actually within
reach, and is honest about the direction that is not. The gap between
those two is real mathematics, not a gap in effort, and closing it is
left where it belongs: open.

\section*{The conjecture, stated precisely}

\begin{conjecture}[Universal continuation generation]
\label{conj:universal}
There exists a finite family of contractions $T_1, \dots, T_n$ on
$X$, fixed once and independent of any particular target, such that
for every nonempty compact $A \subseteq X$ and every $\varepsilon >
0$, some composition $T_{i_k} \circ \cdots \circ T_{i_1}$ applied to
some starting set satisfies $d_H(A, T_{i_k} \circ \cdots \circ
T_{i_1}(\cdot)) < \varepsilon$.
\end{conjecture}

\begin{remark}
This is stated as a conjecture and not as a theorem because it is
one; nothing in this book has proved it, and nothing in this chapter
will. What can be made precise is why it is hard, and that requires
first showing what is \emph{not} hard.
\end{remark}

\section*{The Collage Theorem}

\begin{theorem}[Collage Theorem]
\label{thm:collage}
Let $T_1, \dots, T_n$ be contractions with $\mathcal{H}(A) =
\bigcup_i T_i(A)$ and maximum contraction constant $c < 1$, and let
$A^\ast$ be the attractor of $\mathcal{H}$. Then for any nonempty
compact $A$,
\[
d_H(A, A^\ast) \leq \frac{d_H(A, \mathcal{H}(A))}{1-c}.
\]
\end{theorem}
\begin{proof}
By the triangle inequality and $A^\ast = \mathcal{H}(A^\ast)$,
\[
d_H(A, A^\ast) \leq d_H(A, \mathcal{H}(A)) + d_H(\mathcal{H}(A), \mathcal{H}(A^\ast)).
\]
By Proposition~3 of Chapter 1 ($\mathcal{H}$ is a contraction on
shapes), $d_H(\mathcal{H}(A), \mathcal{H}(A^\ast)) \leq c\, d_H(A,
A^\ast)$. Substituting,
\[
d_H(A, A^\ast) \leq d_H(A, \mathcal{H}(A)) + c\, d_H(A, A^\ast),
\]
and rearranging, since $c < 1$, gives
$(1-c)\, d_H(A,A^\ast) \leq d_H(A, \mathcal{H}(A))$, which is the
claim.
\end{proof}

\begin{remark}
Nothing here was invented for this chapter; the theorem is a direct
three-line consequence of facts Chapter 1 already established, and it
is worth noting how much it buys for how little it costs. It converts
a question about the attractor, an object defined only as a limit, so
generally hard to reason about directly, into a question about a
single application of $\mathcal{H}$ to a set already in hand. Finding
an operator family whose attractor is close to a target $A$ reduces to
finding one whose \emph{one-step collage} $\mathcal{H}(A)$ is already
close to $A$, an ordinary geometric fitting problem rather than a
limiting one.
\end{remark}

\section*{Existence of an approximating family}

\begin{theorem}[Every compact target is approximable]
\label{thm:existence}
For every nonempty compact $A \subseteq X = \mathbb{R}^n$ and every
$\varepsilon > 0$, there exists a finite family of contractions
(depending on $A$ and $\varepsilon$) whose attractor $A^\ast$
satisfies $d_H(A, A^\ast) < \varepsilon$.
\end{theorem}
\begin{proof}[Proof sketch]
Fix $c \in (0,1)$ and let $\{x_1, \dots, x_m\}$ be a finite
$\delta$-net of $A$ for $\delta$ small relative to $\varepsilon(1-c)$,
which exists since $A$ is compact, hence totally bounded. Fix a
bounded universe $U \supseteq A$ and choose contractions $T_i$, each
with ratio $c$ and fixed point $x_i$, scaled so that $T_i(U)$ lies
within a small neighborhood of $x_i$ (achievable by taking the
contraction ratio applied to $U$ small enough relative to $\delta$,
which can be arranged independently for each $i$ while keeping the
overall maximum ratio bounded away from $1$). The resulting collage
$\mathcal{H}(A) = \bigcup_i T_i(A)$ then lies within a bounded
multiple of $\delta$ of the net $\{x_1,\dots,x_m\}$, which is itself
within $\delta$ of $A$, giving $d_H(A, \mathcal{H}(A))$ small relative
to $\varepsilon(1-c)$. Theorem~\ref{thm:collage} then gives $d_H(A,
A^\ast) < \varepsilon$. This is the standard construction underlying
fractal image compression, given in full quantitative detail in
\cite{barnsley1988}, and is not reproduced with complete constants
here.
\end{proof}

\begin{remark}
This is real progress, and it is worth being precise about exactly
how much. It says the space of Hutchinson attractors is dense, in
Hausdorff distance, in the space of nonempty compact subsets of $X$:
nothing worth generating is permanently out of reach. What it does
not say, and what Conjecture~\ref{conj:universal} actually asks for,
is that a \emph{single} family, fixed before any target is known,
suffices for every target. The construction above produces a
different family for each new target and each new tolerance; as
$\varepsilon \to 0$ or as $A$ grows more intricate, the number of
generators $m$ used in the proof generally grows without bound. The
existence theorem answers ``can this be reached at all,'' which
Theorem~\ref{thm:existence} settles completely. The conjecture asks
``can it be reached with tools chosen in advance,'' which is a
categorically stronger demand, and this chapter has not moved that
needle.
\end{remark}

\section*{What kind of gap this is}

The distinction is not a technicality, and it has a name in
neighboring fields worth being explicit about rather than gestured at.
The Stone--Weierstrass theorem guarantees that polynomials are dense
in the continuous functions on a compact interval: a fixed class of
building blocks, closed under the operations that generate it, dense
in everything worth approximating. A universal approximation theorem
for neural networks says something structurally similar: networks
built from a fixed activation function, with enough width, are dense
in the continuous functions. Both are theorems about a fixed generator
class being rich enough for every target, which is exactly the shape
of Conjecture~\ref{conj:universal} and exactly what
Theorem~\ref{thm:existence} does not provide, since its generator
family is chosen after the target is known.

\begin{remark}
Whether an analogue of Stone--Weierstrass or universal approximation
holds for continuation geometries specifically is not something this
chapter's methods can settle, and it should not be assumed to transfer
just because the surface shape of the claim resembles those theorems.
Both classical results rely on structure, algebraic closure under
multiplication for Stone--Weierstrass, compositional depth for neural
networks, that has no established analogue here. Building that
structure, or showing it cannot exist, is what would actually resolve
Conjecture~\ref{conj:universal}, and it remains the open question this
book set out to sharpen rather than close.
\end{remark}

\chapter*{What Remains}
\addcontentsline{toc}{chapter}{What Remains}

This book is unfinished. What follows is a roadmap, not a draft, so
that continuation from here does not require reconstructing the plan
from scratch.

\section*{Chapter 6 (planned): State Insufficiency, Made Explicit}
The sleep/coma motivation used informally since the project's earliest
notes has never been built into an explicit toy system. The chapter
owed is: construct an actual $(X, F, K)$ with two states $x_1, x_2$
(or two histories $h_1, h_2$ landing on the same observable state)
such that $\Gamma(x_1) \neq \Gamma(x_2)$, proved, not narrated. This
is largely already implicit in \emph{Continuation Geometry}, Chapter
5 (the sufficient-statistic material) and in this book's traffic-light
examples, so the honest options are: (a) a short chapter distilling
the sharpest existing instance into a clean standalone theorem, framed
explicitly as ``state insufficiency,'' or (b) a note explaining why a
separate chapter is not needed, since Chapter 5 of the first book
already proves the general case. Decide which before writing anything
new.

\section*{Chapter 7 (planned): Semantic Manifolds}
The most speculative and least-earned material still on the table.
The goal, meaning as a continuation-equivalence class rather than a
point in a fixed space of meanings, requires a specified equivalence
before anything else: candidates include $d_H(\Gamma(h_1),
\Gamma(h_2)) = 0$ (reusing $\sim_\Gamma$ directly, in which case this
chapter may collapse into a corollary of Chapter 2 rather than needing
new machinery) or some notion of isometry or homeomorphism between
continuation spaces (a genuinely stronger and harder-to-work-with
relation, not yet justified over the simpler one). Resist writing a
Theorem with an undefined $\cong$ in it. Settle the equivalence first;
if it turns out to just be $\sim_\Gamma$, say so plainly rather than
introducing new notation to relabel an old idea.

\section*{Chapter 8 (planned): Closing}
A synthesis chapter in the register of \emph{Continuation Geometry}'s
Chapter 11: what was established (Collage Theorem, quotient
invariance, the operator/attractor unification of Chapter 4, the
loss measure of Chapter 3), what was not (the universal generation
conjecture, semantic manifolds if Chapter 7 does not resolve it), and
an honest final assessment of whether ``generative'' turned out to be
the right word for what this book actually built.

\section*{Loose threads not yet assigned to a chapter}
Two things surfaced during drafting and were deliberately not chased:
whether operators named Collapse or Refuse elsewhere in this corpus
are literally contractions in the sense of Chapter 1 (an empirical
question about those specific constructions, not a mathematical one
this book can settle from the armchair); and whether a coordination-
geometry extension of Chapter 4's operator vocabulary, multiple
agents' operators acting jointly on a shared continuation space, is
worth its own chapter or belongs in a third volume alongside
\emph{Continuation Geometry}'s own unfinished coordination-geometry
business. Neither should be started without first deciding whether it
belongs in this book at all.

\newpage
\addcontentsline{toc}{chapter}{References}

\begin{thebibliography}{99}
\bibitem{falconer2003}K.\ Falconer, \emph{Fractal Geometry:
Mathematical Foundations and Applications}, 2nd ed., Wiley, 2003.
\bibitem{barnsley1988}M.\ Barnsley, \emph{Fractals Everywhere},
Academic Press, 1988.
\bibitem{hutchinson1981}J.\ E.\ Hutchinson, ``Fractals and
self-similarity,'' Indiana University Mathematics Journal, 30 (1981).
\end{thebibliography}

\end{document}

